\clearpage
\VerdierPhysicalPageAnchor{0060}%
\VerdierSemanticAnchor{ch1-1-5-14-continuation}
そこで、次のようにおく：

\VerdierSemanticAnchor{ch1-1-5-14-eq3}
\begin{equation}
  \begin{aligned}
    T(0)&=\operatorname{id}_{\mathcal C},\\
    T(n)&=T^n, && n>0,\\
    T(n)&=(T^{-1})^{-n}, && n<0.
  \end{aligned}
  \tag{1.5.14.3}
\end{equation}

\VerdierSemanticAnchor{ch1-1-5-14-cocycle-data}
さらに、整数の各対 $(n,n')$ に対し、同型
\[
  c(n',n)\colon T(n+n')\longrightarrow T(n)T(n')
\]
を次のように定める：

\VerdierSemanticAnchor{ch1-1-5-14-eq4}
\begin{equation}
  \begin{aligned}
    n,n'&\geq 0,
      &c(n',n)&=\operatorname{id},\\
    n,n'&\leq 0,
      &c(n',n)&=\operatorname{id},\\
    n'&\leq 0, n\geq 0, n+n'\leq 0,
      &c(n',n)&=u^{(n)}\mathbin{*}T(n+n'),\\
    n'&\leq 0, n\geq 0, n+n'\geq 0,
      &c(n',n)&=T(n+n')\mathbin{*}u^{(-n')},\\
    n'&\geq 0, n\leq 0, n+n'\leq 0,
      &c(n',n)&=T(n+n')\mathbin{*}v^{(n')},\\
    n'&\geq 0, n\leq 0, n+n'\geq 0,
      &c(n',n)&=v^{(-n)}\mathbin{*}T(n+n').
  \end{aligned}
  \tag{1.5.14.4}
\end{equation}

\VerdierSemanticAnchor{ch1-1-5-14-conclusion}
これによって実際に $\mathbb Z$-圏が定義されたことを確かめるには、
(1.2.1.3) 型の恒等式が成り立つことを検証すればよい。この検証は射
$u^{(n)}$ と $v^{(n)}$ が満たす古典的な恒等式に基づく。詳細は読者に委ねる。

\VerdierSemanticAnchor{ch1-1-6}
\subsection{Koszul 余サイクルと符号規則}

\VerdierSemanticAnchor{ch1-1-6-1}
\noindent\textbf{1.6.1.}\quad
$e_1,\ldots,e_n$ を $\mathbb Z^n$ の基底とする。
$i,j\in[1,n]$、$i\neq j$ に対し、$\mathbb Z^n_{i,j}$ を $e_i$ と
$e_j$ が生成する $\mathbb Z^n$ の部分群とする。制限準同型は準同型

\VerdierSemanticAnchor{ch1-1-6-1-eq1}
\begin{equation}
  H^2(\mathbb Z^n,\mathbb Z/2\mathbb Z)
  \longrightarrow
  \bigoplus_{\substack{i,j\in[1,n]\\ i<j}}
    H^2(\mathbb Z^n_{i,j},\mathbb Z/2\mathbb Z)
  \tag{1.6.1.1}
\end{equation}
を定める。

（標準的な）同型

\VerdierSemanticAnchor{ch1-1-6-1-eq2}
\begin{equation}
  H^2(\mathbb Z^n_{i,j},\mathbb Z/2\mathbb Z)
  \simeq \mathbb Z/2\mathbb Z
  \tag{1.6.1.2}
\end{equation}
がある。

\VerdierSemanticAnchor{ch1-1-6-1-conclusion}
準同型 (1.6.1.1) は同型である。
